\Chapter{52}

\Verse{1}
{
    \zA{话说贾母道：“正是这个了。}
}
{
    \eA{But let us return to our story. "Quite so!" was the reply with which dowager
        lady Chia (greeted lady Feng's proposal).}
}
{
}

\Verse{2}
{
    \zA{上次我要说这话，我见你们大事多，如今又添出 些事来，你们固然不敢抱怨，未免想着我只顾疼这些小孙子孙女儿们，就不体贴你 们这当家人了。}
}
{
    \eA{"I meant the other day to have suggested this arrangement, but I saw that
        every one of you had so many urgent matters to attend to, (and I thought)
        that although you would not presume to bear me a grudge, were several duties
        now again superadded, you would unavoidably imagine that I only regarded
        those young grandsons and granddaughters of mine, and had no consideration
        for any of you,}
}
{
}

\Verse{3}
{
    \zA{你既这么说出来，便好了。”}
}
{
    \eA{who have to look after the house. But since you make this suggestion
        yourself,}
}
{
}

\Verse{4}
{
    \zA{因此时薛姨妈李婶娘都在座，邢夫人 及尤氏等也都过来请安，还未过去，贾母因向王夫人等说道：“今日我才说这话，
        素日我不说：一则怕逞了凤丫头的脸，二则众人不服。}
}
{
    \eA{it's all right." And seeing that Mrs. Hsüeh, and 'sister-in-law' Li were
        sitting with her, and that Madame Hsing, and Mrs. Yu and the other ladies,
        who had also crossed over to pay their respects, had not as yet gone to
        their quarters, old lady Chia broached the subject with Madame Wang,}
}
{
}

\Verse{5}
{
    \zA{今日你们都在这里，都是经 过妯娌姑嫂的，还有他这么想得到的没有？”}
}
{
    \eA{and the rest of the company. "I've never before ventured to give utterance
        to the remarks that just fell from my lips," she said, "as first of all I
        was in fear and trembling lest I should have made that girl Feng more
        presumptuous than ever,}
}
{
}

\Verse{6}
{
    \zA{薛姨妈、李婶娘、尤氏齐笑说：“真 个少有!别人不过是礼上的面情儿，实在他是真疼小姑子小叔子。}
}
{
    \eA{and next, lest I should have incurred the displeasure of one and all of you.
        But since you're all here to-day, and every one of you knows what brothers'
        wives and husbands' sisters mean, is there (I ask) any one besides her as
        full of forethought?"}
}
{
}

\Verse{7}
{
    \zA{就是老太太跟前， 也是真孝顺。”}
}
{
    \eA{Mrs. Hsüeh, 'sister-in-law' Li and Mrs. Yu smiled with one consent. "There
        are indeed but few like her!"}
}
{
}

\Verse{8}
{
    \zA{贾母点头叹道：“我虽疼他，我又怕他太伶俐了，也不是好事。”}
}
{
    \eA{they cried. "That of others is simply a conventional 'face' affection, but
        she is really fond of her husband's sisters and his young brother. In fact,
        she's as genuinely filial with you,}
}
{
}

\Verse{9}
{
    \zA{凤姐儿忙笑道：“这话老祖宗说差了。}
}
{
    \eA{venerable senior." Dowager lady Chia nodded her head. "Albeit I'm fond of
        her," she sighed, "I can't, on the other hand,}
}
{
}

\Verse{10}
{
    \zA{世人都说：‘太伶俐聪明怕活不长。’}
}
{
    \eA{help distrusting that excessive shrewdness of hers, for it isn't a good
        thing." "You're wrong there, worthy ancestor,"}
}
{
}

\Verse{11}
{
    \zA{世人 都说，世人都信，独老祖宗不当说，不当信。}
}
{
    \eA{lady Feng laughed with alacrity. "People in the world as a rule maintain
        that 'too shrewd and clever a person can't, it is feared, live long.'}
}
{
}

\Verse{12}
{
    \zA{老祖宗只有伶俐聪明过我十倍的，怎 么如今这么福寿双全的？}
}
{
    \eA{Now what people of the world invariably say people of the world invariably
        believe. But of you alone, my dear senior, can no such thing be averred or
        believed.}
}
{
}

\Verse{13}
{
    \zA{只怕我明儿还胜老祖宗一倍呢。}
}
{
    \eA{For there you are, ancestor mine, a hundred times sharper and cleverer than
        I;}
}
{
}

\Verse{14}
{
    \zA{我活一千岁后，等老祖宗 归了西，我才死呢。”}
}
{
    \eA{and how is it that you now enjoy both perfect happiness and longevity? But I
        presume that I shall by and bye excel you by a hundredfold,}
}
{
}

\Verse{15}
{
    \zA{贾母笑道：“众人都死了，单剩咱们两个老妖精，有什么意 思！”}
}
{
    \eA{and die at length, after a life of a thousand years, when you venerable
        senior shall have departed from these mortal scenes!" "After every one is
        dead and gone," dowager lady Chia laughingly observed,}
}
{
}

\Verse{16}
{
    \zA{说的众人都笑了。}
}
{
    \eA{"what pleasure will there be, if two antiquated elves,}
}
{
}

\Verse{17}
{
    \zA{宝玉因惦记着晴雯等事，便先回园里来。}
}
{
    \eA{like you and I will be, remain behind?" This joke excited general mirth. But
        so concerned was Pao-yü about Ch'ing Wen and other matters that he was the
        first to make a move and return into the garden. On his arrival at his
        quarters,}
}
{
}

\Verse{18}
{
    \zA{到了屋中，药香满室，一人不见，只 有晴雯独卧于炕上，脸上烧的飞红。}
}
{
    \eA{he found the rooms full of the fragrance emitted by the medicines. Not a
        soul did he, however, see about. Ch'ing Wen was reclining all alone on the
        stove-couch. Her face was feverish and red.}
}
{
}

\Verse{19}
{
    \zA{又摸了一摸，只觉烫手，忙又向炉上将手烘暖， 伸进被去摸了一摸身上，也是火热。}
}
{
    \eA{When he came to touch it, his hand experienced a scorching sensation.
        Retracing his steps therefore towards the stove, he warmed his hands and
        inserted them under the coverlet and felt her. Her body as well was as hot
        as fire.}
}
{
}

\Verse{20}
{
    \zA{因说道：“别人去了也罢，麝月秋纹也这么无 情，各自去了？”}
}
{
    \eA{"If the others have left," he then remarked, "there's nothing strange about
        it, but are She Yüeh and Ch'iu Wen too so utterly devoid of feeling as to
        have each gone after her own business?"}
}
{
}

\Verse{21}
{
    \zA{晴雯道：“秋纹是我撵了他去吃饭了，麝月是方才平儿来找他出 去了，两个人鬼鬼祟祟的，不知说什么。}
}
{
    \eA{"As regards Ch'iu Wen," Ch'ing Wen explained, "I told her to go and have her
        meal. And as for She Yüeh, P'ing Erh came just now and called her out of
        doors and there they are outside confabbing in a mysterious way! What the
        drift of their conversation can be I don't know. But they must be talking
        about my having fallen ill,}
}
{
}

\Verse{22}
{
    \zA{必是说我病了不出去。”}
}
{
    \eA{and my not leaving this place to go home." "P'ing Erh isn't that sort of
        person,"}
}
{
}

\Verse{23}
{
    \zA{宝玉道：“平儿 不是那样人。}
}
{
    \eA{Pao-yü pleaded. "Besides, she had no idea whatever about your illness,}
}
{
}

\Verse{24}
{
    \zA{况且他并不知你病特来瞧你，想来一定是找麝月来说话，偶然见你病 了，随口说特瞧你的病，这也是人情乖觉取和儿的常事。}
}
{
    \eA{so that she couldn't have come specially to see how you were getting on. I
        fancy her object was to look up She Yüeh to hobnob with her, but finding
        unexpectedly that you were not up to the mark, she readily said that she had
        come on purpose to find what progress you were making. This was quite a
        natural thing for a person with so wily a disposition to say,}
}
{
}

\Verse{25}
{
    \zA{便不出去，有不是，与他 何干？}
}
{
    \eA{for the sake of preserving harmony. But if you don't go home, it's none of
        her business.}
}
{
}

\Verse{26}
{
    \zA{你们素日又好，断不肯为这无干的事伤和气。”}
}
{
    \eA{You two have all along been, irrespective of other things, on such good
        terms that she could by no means entertain any desire to injure the friendly
        relations which exist between you,}
}
{
}

\Verse{27}
{
    \zA{晴雯道：“这话也是，只是 疑他为什么忽然又瞒起我来？”}
}
{
    \eA{all on account of something that doesn't concern her." "Your remarks are
        right enough," Ch'ing Wen rejoined, "but I do suspect her, as why did she
        too start,}
}
{
}

\Verse{28}
{
    \zA{宝玉笑道：“等我从后门出去，到那窗户根下听听 说些什么，来告诉你。”}
}
{
    \eA{all of a sudden, imposing upon me?" "Wait, I'll walk out by the back door,"
        Pao-yü smiled, "and go to the foot of the window, and listen to what she's
        saying. I'll then come and tell you."}
}
{
}

\Verse{29}
{
    \zA{说着，果从后门出去至窗下，潜听麝月悄悄问道：“你怎么就得了的？”}
}
{
    \eA{Speaking the while, he, in point of fact, sauntered out of the back door;
        and getting below the window, he lent an ear to their confidences. "How did
        you manage to get it?"}
}
{
}

\Verse{30}
{
    \zA{平儿 道：“那日彼时洗手时不见了，二奶奶就不许吵嚷；出了园子，即刻就传给园里各 处的妈妈们，小心访查。}
}
{
    \eA{She Yueh inquired with gentle voice. "When I lost sight of it on that day
        that I washed my hands," P'ing Erh answered, "our lady Secunda wouldn't let
        us make a fuss. But the moment she left the garden, she there and then sent
        word to the nurses, stationed in the various places,}
}
{
}

\Verse{31}
{
    \zA{我们只疑惑邢姑娘的丫头，本来又穷，只怕小孩子家没见 过，拿起来是有的，再不料定是你们这里的。}
}
{
    \eA{to institute careful search. Our suspicions, however, fell upon Miss Hsing's
        maid, who has ever also been poverty-stricken; surmising that a young girl
        of her age, who had never set eyes upon anything of the kind, may possibly
        have picked it up and taken it.}
}
{
}

\Verse{32}
{
    \zA{幸而二奶奶没有在屋里，你们这里的 宋妈去了，拿着这支镯子，说是小丫头坠儿偷起来的，被他看见，来回二奶奶的。}
}
{
    \eA{But never did we positively believe that it could be some one from this
        place of yours! Happily, our lady Secunda wasn't in the room, when that
        nurse Sung who is with you here went over, and said, producing the bracelet,
        'that the young maid, Chui Erh, had stolen it, and that she had detected
        her,}
}
{
}

\Verse{33}
{
    \zA{我赶忙接了镯子。}
}
{
    \eA{and come to lay the matter before our lady Secunda.}
}
{
}

\Verse{34}
{
    \zA{想了一想：宝玉是偏在你们身上留心用意、争胜要强的，那一年 有个良儿偷玉，刚冷了这二年，闲时还常有人提起来趁愿；这会子又跑出一个偷金
        子的来了，而且更偷到街坊家去了!偏是他这么着，偏是他的人打嘴。}
}
{
    \eA{I promptly took over the bracelet from her; and recollecting how imperious
        and exacting Pao-yü is inclined to be, fond and devoted as he is to each and
        all of you; how the jade which was prigged the other year by a certain Liang
        Erh, is still, just as the matter has cooled down for the last couple of
        years, canvassed at times by some people eager to serve their own ends; how
        some one has now again turned up to purloin this gold trinket;}
}
{
}

\Verse{35}
{
    \zA{所以我倒忙 叮咛宋妈千万别告诉宝玉，只当没有这事，总别和一个人提起。}
}
{
    \eA{how it was filched, to make matters worse, from a neighbour's house; how as
        luck would have it, she took this of all things; and how it happened to be
        his own servant to give him a slap on his mouth,}
}
{
}

\Verse{36}
{
    \zA{第二件，老太太、 太太听了生气。}
}
{
    \eA{I hastened to enjoin nurse Sung to, on no account whatever, let Pao-yü know
        anything about it, but simply pretend that nothing of the kind had
        transpired,}
}
{
}

\Verse{37}
{
    \zA{三则袭人和你们也不好看。}
}
{
    \eA{and to make no mention of it to any single soul. In the second place,' (I
        said),}
}
{
}

\Verse{38}
{
    \zA{所以我回二奶奶只说：‘我往大奶奶那 里去来着，谁知镯子褪了口，丢在草根底下，雪深了没看见。}
}
{
    \eA{'our dowager lady and Madame Wang would get angry, if they came to hear
        anything. Thirdly, Hsi Jen as well as yourselves would not also cut a very
        good figure.' Hence it was that in telling our lady Secunda, I merely
        explained 'that on my way to our senior mistress,'}
}
{
}

\Verse{39}
{
    \zA{今儿雪化尽了，黄澄 澄的映着日头，还在那里呢，我就拣了起来。’}
}
{
    \eA{the bracelet got unclasped, without my knowing it; that it fell among the
        roots of the grass; that there was no chance of seeing it while the snow was
        deep,}
}
{
}

\Verse{40}
{
    \zA{二奶奶也就信了，所以我来告诉你 们。}
}
{
    \eA{but that when the snow completely disappeared to-day there it glistened, so
        yellow and bright, in the rays of the sun,}
}
{
}

\Verse{41}
{
    \zA{你们以后防着他些，别使唤他到别处去。}
}
{
    \eA{in precisely the very place where it had dropped, and that I then picked it
        up.' Our lady Secunda at once credited my version.}
}
{
}

\Verse{42}
{
    \zA{等袭人回来，你们商议着，变个法子 打发出去就完了。”}
}
{
    \eA{So here I come to let you all know so as to be henceforward a little on your
        guard with her, and not get her a job anywhere else. Wait until Hsi Jen's
        return,}
}
{
}

\Verse{43}
{
    \zA{麝月道：“这小娼妇也见过些东西，怎么这么眼浅？”}
}
{
    \eA{and then devise means to pack her off, and finish with her." "This young
        vixen has seen things of this kind before,"}
}
{
}

\Verse{44}
{
    \zA{平儿道： “究竟这镯子能多重!原是二奶奶的，说这叫做‘虾须镯’，倒是这颗珠子重了。}
}
{
    \eA{She Yüeh ejaculated, "and how is it that she was so shallow-eyed?" "What
        could, after all, be the weight of this bracelet?" P'ing Erh observed. "It
        was once our lady Secunda's. She says that this is called the
        'shrimp-feeler'-bracelet.}
}
{
}

\Verse{45}
{
    \zA{晴雯那蹄子是块爆炭，要告诉了他，他是忍不住的，一时气上来，或打或骂，依旧 嚷出来，所以单告诉你留心就是了。”}
}
{
    \eA{But it's the pearl, which increases its weight. That minx Ch'ing Wen is as
        fiery as a piece of crackling charcoal, so were anything to be told her, she
        may, so little able is she to curb her temper, flare up suddenly into a
        huff, and beat or scold her, and kick up as much fuss as she ever has done
        before.}
}
{
}

\Verse{46}
{
    \zA{说着，便作辞而去。}
}
{
    \eA{That's why I simply tell you. Exercise due care, and it will be all right."}
}
{
}

\Verse{47}
{
    \zA{宝玉听了，又喜又气又叹：喜的是平儿竟能体帖自己的心；气的是坠儿小窃； 叹的是坠儿那样伶俐，做出这丑事来。}
}
{
    \eA{With this warning, she bid her farewell and went on her way. Her words
        delighted, vexed and grieved Pao-yü. He felt delighted, on account of the
        consideration shown by P'ing Erh for his own feelings. Vexed, because Chui
        Erh had turned out a petty thief. Grieved, that Chui Erh, who was otherwise
        such a smart girl,}
}
{
}

\Verse{48}
{
    \zA{因而回至房中，把平儿之话一长一短告诉了 晴雯，又说：“他说你是个要强的，如今病了，听了这话，越发要添病的，等好了 再告诉你。”}
}
{
    \eA{should have gone in for this disgraceful affair. Returning consequently into
        the house, he told Ch'ing Wen every word that P'ing Erh had uttered. "She
        says," he went on to add, "that you're so fond of having things all your own
        way that were you to hear anything of this business, now that you are ill,
        you would get worse,}
}
{
}

\Verse{49}
{
    \zA{晴雯听了，果然气的蛾眉倒蹙，凤眼圆睁，即时就叫坠儿。}
}
{
    \eA{and that she only means to broach the subject with you, when you get quite
        yourself again." Upon hearing this, Ch'ing Wen's ire was actually stirred
        up,}
}
{
}

\Verse{50}
{
    \zA{宝玉忙劝 道：“这一喊出来，岂不辜负了平儿待你我的心呢？}
}
{
    \eA{and her beautiful moth-like eyebrows contracted, and her lovely phoenix eyes
        stared wide like two balls. So she immediately shouted out for Chui Erh. "If
        you go on bawling like that,"}
}
{
}

\Verse{51}
{
    \zA{不如领他这个情，过后打发他 出去就完了。”}
}
{
    \eA{Pao-yü hastily remonstrated with her, "won't you show yourself ungrateful
        for the regard with which P'ing Erh has dealt with you and me?}
}
{
}

\Verse{52}
{
    \zA{晴雯道：“虽如此说，只是这气如何忍得住？”}
}
{
    \eA{Better for us to show ourselves sensible of her kindness and by and bye pack
        the girl off, and finish." "Your suggestion is all very good,"}
}
{
}

\Verse{53}
{
    \zA{宝玉道：“这有什 么气的？}
}
{
    \eA{Ch'ing Wen demurred, "but how could I suppress this resentment?"}
}
{
}

\Verse{54}
{
    \zA{你只养病就是了。”}
}
{
    \eA{"What's there to feel resentment about?"}
}
{
}

\Verse{55}
{
    \zA{晴雯服了药，至晚间又服了二和，夜间虽有些汗，还未见效，仍是发烧头疼鼻 塞声重。}
}
{
    \eA{Pao-yü asked. "Just you take good care of yourself; it's the best thing you
        can do." Ch'ing Wen then took her medicine. When evening came, she had
        another couple of doses. But though in the course of the night,}
}
{
}

\Verse{56}
{
    \zA{次日，王太医又来诊视，另加减汤剂。}
}
{
    \eA{she broke out into a slight perspiration, she did not see any change for the
        better in her state.}
}
{
}

\Verse{57}
{
    \zA{虽然稍减了烧，仍是头疼。}
}
{
    \eA{Still she felt feverish, her head sore, her nose stopped, her voice hoarse.
        The next day,}
}
{
}

\Verse{58}
{
    \zA{宝玉便 命麝月取鼻烟来：“给他闻些，痛打几个嚏喷就通快了。”}
}
{
    \eA{Dr. Wang came again to examine her pulse and see how she was getting on.
        Besides other things, he increased the proportions of certain medicines in
        the decoction and reduced others;}
}
{
}

\Verse{59}
{
    \zA{麝月果真去取了一个金 镶双金星玻璃小扁盒儿来递给宝玉。}
}
{
    \eA{but in spite of her fever having been somewhat brought down, her head
        continued to ache as much as ever. "Go and fetch the snuff," Pao-yü said to
        She Yüeh,}
}
{
}

\Verse{60}
{
    \zA{宝玉便揭开盒盖，里面是个西洋珐琅的黄发赤 身女子，两肋又有肉翅，里面盛着些真正上等洋烟。}
}
{
    \eA{"and give it to her to sniff. She'll feel more at ease after she has had
        several strong sneezes." She Yüeh went, in fact, and brought a flat crystal
        bottle, inlaid with a couple of golden stars, and handed it to Pao-yü.
        Pao-yü speedily raised the cover of the bottle.}
}
{
}

\Verse{61}
{
    \zA{晴雯只顾看画儿，宝玉道：“闻 些，走了气就不好了。”}
}
{
    \eA{Inside it, he discovered, represented on western enamel, a fair-haired young
        girl, in a state of nature, on whose two sides figured wings of flesh.}
}
{
}

\Verse{62}
{
    \zA{睛雯听说，忙用指甲挑了些抽入鼻中。}
}
{
    \eA{This bottle contained some really first-rate foreign snuff. Ch'ing Wen's
        attention was fixedly concentrated on the representation.}
}
{
}

\Verse{63}
{
    \zA{不见怎么，便又多 多挑了些抽入。}
}
{
    \eA{"Sniff a little!" Pao-yü urged. "If the smell evaporates, it won't be worth
        anything."}
}
{
}

\Verse{64}
{
    \zA{忽觉鼻中一股酸辣，透入囟门，接连打了五六个嚏喷，眼泪鼻涕登 时齐流。}
}
{
    \eA{Ch'ing Wen, at his advice, promptly dug out a little with her nail, and
        applied it to her nose. But with no effect. So digging out again a good
        quantity of it, she pressed it into her nostrils.}
}
{
}

\Verse{65}
{
    \zA{晴雯忙收了盒子，笑道：“了不得，辣!快拿纸来。”}
}
{
    \eA{Then suddenly she experienced a sensation in her nose as if some pungent
        matter had penetrated into the very duct leading into the head,}
}
{
}

\Verse{66}
{
    \zA{早有小丫头子递过 一搭子细纸，晴雯便一张一张的拿来醒鼻子。}
}
{
    \eA{and she sneezed five or six consecutive times, until tears rolled down from
        her eyes and mucus trickled from her nostrils. Ch'ing Wen hastily put the
        bottle away.}
}
{
}

\Verse{67}
{
    \zA{宝玉笑问：“如何？”}
}
{
    \eA{"It's dreadfully pungent!" she laughed. "Bring me some paper,}
}
{
}

\Verse{68}
{
    \zA{晴雯笑道：“果 然通快些。}
}
{
    \eA{quick!" A servant-girl at once handed her a pile of fine paper.}
}
{
}

\Verse{69}
{
    \zA{只是太阳还疼。”}
}
{
    \eA{Ch'ing Wen extracted sheet after sheet,}
}
{
}

\Verse{70}
{
    \zA{宝玉笑道：“越发尽用西洋药治一治，只怕就好了。”}
}
{
    \eA{and blew her nose. "Well," said Pao-yü smiling, "how are you feeling now?"
        "I'm really considerably relieved." Ch'ing Wen rejoined laughing.}
}
{
}

\Verse{71}
{
    \zA{说着，便命麝月：“往二奶奶要去，就说我说了，姐姐那里常有那西洋贴头疼的膏 子药，叫做‘依佛哪’，找寻一点儿。”}
}
{
    \eA{"The only thing is that my temples still hurt me." "Were you to treat
        yourself exclusively with western medicines, I'm sure you'd get all right,"
        Pao-yü added smilingly. Saying this,}
}
{
}

\Verse{72}
{
    \zA{麝月答应去了，半日，果然拿了半节来。}
}
{
    \eA{"Go," he accordingly desired She Yüeh, "to our lady Secunda, and ask her for
        some. Tell her that I spoke to you about them.}
}
{
}

\Verse{73}
{
    \zA{便去找了一块红缎子角儿，铰了两块指顶大的圆式，将那药烤和了，用簪挺摊上。}
}
{
    \eA{My cousin over there often uses some western plaster, which she applies to
        her temples when she's got a headache. It's called 'I-fo-na.' So try and get
        some of it!" She Yüeh expressed her readiness. After a protracted absence,}
}
{
}

\Verse{74}
{
    \zA{晴雯自拿着一面靶儿镜子贴在两太阳上。}
}
{
    \eA{she, in very deed, came back with a small bit of the medicine; and going
        quickly for a piece of red silk cutting,}
}
{
}

\Verse{75}
{
    \zA{麝月笑道：“病的蓬头鬼一样，如今贴了 这个，倒俏皮了!二奶奶贴惯了，倒不大显。”}
}
{
    \eA{she got the scissors and slit two round slips off as big as the tip of a
        finger. After which, she took the medicine, and softening it by the fire,
        she spread it on them with a hairpin.}
}
{
}

\Verse{76}
{
    \zA{说毕，又问宝玉道：“二奶奶说了： 明儿是舅老爷的生日，太太说了叫你去呢。}
}
{
    \eA{Ch'ing Wen herself laid hold of a looking-glass with a handle and stuck the
        bits on both her temples. "While you were lying sick," She Yüeh laughed,
        "you looked like a mangy-headed devil!}
}
{
}

\Verse{77}
{
    \zA{明儿穿什么衣裳？}
}
{
    \eA{But with this stuff on now you present a fine sight!}
}
{
}

\Verse{78}
{
    \zA{今儿晚上好打点齐备 了，省的明儿早起费手。”}
}
{
    \eA{As for our lady Secunda she has been so much in the habit of sticking these
        things about her that they don't very much show off with her!"}
}
{
}

\Verse{79}
{
    \zA{宝玉道：“什么顺手就是什么罢了。}
}
{
    \eA{This joke over, "Our lady Secunda said," she resumed, addressing herself to
        Pao-yü,}
}
{
}

\Verse{80}
{
    \zA{一年闹生日也闹不 清。”}
}
{
    \eA{"'that to-morrow is your maternal uncle's birthday, and that our mistress,
        your mother,}
}
{
}

\Verse{81}
{
    \zA{说着，便起身出房，往惜春屋里去看画儿。}
}
{
    \eA{asked her to tell you to go over. That whatever clothes you will put on
        to-morrow should be got ready to-night,}
}
{
}

\Verse{82}
{
    \zA{刚到院门外边，忽见宝琴小丫头名小螺的从那边过去。}
}
{
    \eA{so as to avoid any trouble in the morning.'" "Anything that comes first to
        hand," Pao-yü observed, "will do well enough!}
}
{
}

\Verse{83}
{
    \zA{宝玉忙赶上问：“那里 去？”}
}
{
    \eA{There's no getting, the whole year round, at the end of all the fuss of
        birthdays!"}
}
{
}

\Verse{84}
{
    \zA{小螺笑道：“我们二位姑娘都在林姑娘屋里呢，我如今也往那里去。”}
}
{
    \eA{Speaking the while, he rose to his feet and left the room with the idea of
        repairing to Hsi Ch'un's quarters to have a look at the painting. As soon as
        he got outside the door of the court-yard,}
}
{
}

\Verse{85}
{
    \zA{宝玉 听了，转步也便和他往潇湘馆来。}
}
{
    \eA{he unexpectedly spied Pao-ch'in's young maid, Hsiao Lo by name, crossing
        over from the opposite direction.}
}
{
}

\Verse{86}
{
    \zA{不但宝钗姐妹在此，且连岫烟也在那里。}
}
{
    \eA{Pao-yü, with rapid step, strode up to her, and inquired of her whither she
        was going. "Our two young ladies," Hsiao Lo answered with a smile,}
}
{
}

\Verse{87}
{
    \zA{四人团 坐在熏笼上叙家常。}
}
{
    \eA{"are in Miss Lin's rooms; so I'm also now on my way thither."}
}
{
}

\Verse{88}
{
    \zA{紫鹃倒坐在暖阁里，临窗户做针线。}
}
{
    \eA{Catching this answer, Pao-yü wheeled round and came at once with her to the
        Hsiao Hsiang Lodge. Here not only did he find Pao-ch'ai and her cousin,}
}
{
}

\Verse{89}
{
    \zA{一见他来，都笑说：“又 来了一个!没了你的坐处了。”}
}
{
    \eA{but Hsing Chou-yen as well. The quartet was seated in a circle on the
        warming-frame; carrying on a friendly chat on everyday domestic matters;}
}
{
}

\Verse{90}
{
    \zA{宝玉笑道：“好一幅‘冬闺集艳图’!可惜我迟来了。}
}
{
    \eA{while Tzu Chüan was sitting in the winter apartment, working at some
        needlework by the side of the window. The moment they caught a glimpse of
        him,}
}
{
}

\Verse{91}
{
    \zA{横竖这屋子比各屋子暖，这椅子坐着并不冷。”}
}
{
    \eA{their faces beamed with smiles. "There comes some one else!" they cried.
        "There's no room for you to sit!"}
}
{
}

\Verse{92}
{
    \zA{说着，便坐在黛玉常坐的地方，上 搭着灰鼠椅搭一张椅上。}
}
{
    \eA{"What a fine picture of beautiful girls, in the winter chamber!" Pao-yü
        smiled. "It's a pity I come a trifle too late! This room is, at all events,}
}
{
}

\Verse{93}
{
    \zA{因见暖阁之中有一玉石条盆，里面攒三聚五栽着一盆单瓣 水仙，宝玉便极口赞道：“好花!这屋子越暖，这花香的越浓。}
}
{
    \eA{so much warmer than any other, that I won't feel cold if I plant myself on
        this chair." So saying, he made himself comfortable on a favourite chair of
        Tai-yü's over which was thrown a grey squirrel cover. But noticing in the
        winter apartment a jadestone bowl,}
}
{
}

\Verse{94}
{
    \zA{怎么昨儿没见？”}
}
{
    \eA{full of single narcissi, in clusters of three or five,}
}
{
}

\Verse{95}
{
    \zA{黛玉笑道：“这是你家的大总管赖大奶奶送薛二姑娘的两盆水仙、两盆腊梅：他送 了我一盆水仙，送了云丫头一盆腊梅。}
}
{
    \eA{Pao-yü began praising their beauty with all the language he could command.
        "What lovely flowers!" he exclaimed. "The warmer the room gets, the stronger
        is the fragrance emitted by these flowers! How is it I never saw them
        yesterday?"}
}
{
}

\Verse{96}
{
    \zA{我原不要的，又恐辜负了他的心。}
}
{
    \eA{"These are," Tai-yü laughingly explained, "from the two pots of narcissi,
        and two pots of allspice,}
}
{
}

\Verse{97}
{
    \zA{你若要， 我转送你如何？”}
}
{
    \eA{sent to Miss Hsüeh Secunda by the wife of Lai Ta, the head butler in your
        household.}
}
{
}

\Verse{98}
{
    \zA{宝玉道：“我屋里却有两盆，只是不及这个。}
}
{
    \eA{Of these, she gave me a pot of narcissi; and to that girl Yün, a pot of
        allspice. I didn't at first mean to keep them,}
}
{
}

\Verse{99}
{
    \zA{琴妹妹送你的，如 何又转送人，这个断断使不得。”}
}
{
    \eA{but I was afraid of showing no consideration for her kind attention. But if
        you want them, I'll, in my turn, present them to you.}
}
{
}

\Verse{100}
{
    \zA{黛玉道：“我一日药铞子不离火，我竟是药培着 呢，哪里还搁的住花香来熏？}
}
{
    \eA{Will you have them; eh?" "I've got two pots of them in my rooms," Pao-yü
        replied, "but they're not up to these. How is it you're ready to let others
        have what cousin Ch'in has given you? This can on no account do!"}
}
{
}

\Verse{101}
{
    \zA{越发弱了。}
}
{
    \eA{"With me here," Tai-yü added,}
}
{
}

\Verse{102}
{
    \zA{况且这屋子里一股药香，反把这花香搅坏 了。}
}
{
    \eA{"the medicine pot never leaves the fire, the whole day long. I'm only kept
        together by medicines. So how could I ever stand the smell of flowers
        bunging my nose?}
}
{
}

\Verse{103}
{
    \zA{不如你抬了去，这花儿倒清净了，没什么杂味来搅他。”}
}
{
    \eA{It makes me weaker than ever. Besides, if there's the least whiff of
        medicines in this room, it will, contrariwise, spoil the fragrance of these
        flowers.}
}
{
}

\Verse{104}
{
    \zA{宝玉笑道：“我屋里 今儿也有个病人煎药呢。}
}
{
    \eA{So isn't it better that you should have them carried away? These flowers
        will then breathe a purer atmosphere,}
}
{
}

\Verse{105}
{
    \zA{你怎么知道的？”}
}
{
    \eA{and won't have any mixture of smells to annoy them."}
}
{
}

\Verse{106}
{
    \zA{黛玉笑道：“这说奇了。}
}
{
    \eA{"I've also got now some one ill in my place," Pao-yü retorted with a smile,}
}
{
}

\Verse{107}
{
    \zA{我原是无心话： 谁知你屋里的事？}
}
{
    \eA{"and medicines are being decocted. How comes it you happen to know nothing
        about it?"}
}
{
}

\Verse{108}
{
    \zA{你不早来听古记儿，这会子来了，自惊自怪的。”}
}
{
    \eA{"This is strange!" Tai-yü laughed. "I was really speaking quite
        thoughtlessly; for who ever knows what's going on in your apartments?}
}
{
}

\Verse{109}
{
    \zA{宝玉笑道：“咱们明儿下一社又有了题目了：就咏水仙、腊梅。”}
}
{
    \eA{But why do you, instead of getting here a little earlier to listen to old
        stories, come at this moment to bring trouble and vexation upon your own
        self?" Pao-yü gave a laugh. "Let's have a meeting to-morrow,"}
}
{
}

\Verse{110}
{
    \zA{黛玉听了， 笑道：“罢，罢!再不敢做诗了。}
}
{
    \eA{he proposed, "for we've also got the themes. Let's sing the narcissus and
        allspice." "Never mind, drop that!"}
}
{
}

\Verse{111}
{
    \zA{做一回，罚一回，没的怪羞的。”}
}
{
    \eA{Tai-yü rejoined, upon hearing his proposal. "I can't venture to write any
        more verses.}
}
{
}

\Verse{112}
{
    \zA{说着，便两手 握起脸来。}
}
{
    \eA{Whenever I indite any, I'm mulcted. So I'd rather not be put to any great
        shame."}
}
{
}

\Verse{113}
{
    \zA{宝玉笑道：“何苦来，又打趣我做什么？}
}
{
    \eA{While uttering these words she screened her face with both hands. "What's
        the matter?" Pao-yü smiled.}
}
{
}

\Verse{114}
{
    \zA{我还不怕臊呢，你倒握起脸来 了。”}
}
{
    \eA{"Why are you again making fun of me? I'm not afraid of any shame, but, lo,
        you screen your face."}
}
{
}

\Verse{115}
{
    \zA{宝钗因笑道：“下次我邀一社，四个诗题，四个词题。}
}
{
    \eA{"The next time," Pao-ch'ai felt impelled to interpose laughingly, "I convene
        a meeting, we'll have four themes for odes and four for songs;}
}
{
}

\Verse{116}
{
    \zA{每人四首诗，四首词。}
}
{
    \eA{and each one of us will have to write four odes and four roundelays.}
}
{
}

\Verse{117}
{
    \zA{头一个诗题《咏太极图》，限‘一先’的韵，五言排律；要把‘一先’的韵都用尽 了，一个不许剩。”}
}
{
    \eA{The theme of the first ode will treat of the plan of the great extreme; the
        rhyme fixed being 'hsien,' (first), and the metre consisting of five words
        in each line. We'll have to exhaust every one of the rhymes under 'hsien,'
        and mind, not a single one may be left out."}
}
{
}

\Verse{118}
{
    \zA{宝琴笑道：“这一说，可知是姐姐不是真心起社了，这分明是 难人。}
}
{
    \eA{"From what you say," Pao-ch'in smilingly observed, "it's evident that you're
        not in earnest, cousin, in setting the club on foot. It's clear enough that
        your object is to embarrass people.}
}
{
}

\Verse{119}
{
    \zA{要论起来，也强扭的出来，不过颠来倒去，弄些《易经》上的话生填，究竟 有何趣味。}
}
{
    \eA{But as far as the verses go, we could forcibly turn out a few, just by
        higgledy-piggledy taking several passages from the 'Canon of Changes,' and
        inserting them in our own; but, after all, what fun will there be in that
        sort of thing?}
}
{
}

\Verse{120}
{
    \zA{我八岁的时节，跟我父亲到西海沿上买洋货。}
}
{
    \eA{When I was eight years of age, I went with my father to the western seaboard
        to purchase foreign goods. Who'd have thought it, we came across a girl from
        the 'Chen Chen' kingdom. She was in her eighteenth year,}
}
{
}

\Verse{121}
{
    \zA{谁知有个真真国的女孩子， 才十五岁，那脸面就和那西洋画上的美人一样，也披着黄头发，打着联垂，满头带
        着都是玛瑙、珊瑚、猫儿眼、祖母绿，身上穿着金丝织的锁子甲，洋锦袄袖，带着 倭刀也是镶金嵌宝的。}
}
{
    \eA{and her features were just like those of the beauties one sees represented
        in foreign pictures. She had also yellow hair, hanging down, and arranged in
        endless plaits. Her whole head was ornamented with one mass of cornelian
        beads, amber, cats' eyes, and 'grandmother-green-stone.' On her person, she
        wore a chain armour plaited with gold,}
}
{
}

\Verse{122}
{
    \zA{实在画儿上也没他那么好看。}
}
{
    \eA{and a coat, which was up to the very sleeves, embroidered in foreign style.}
}
{
}

\Verse{123}
{
    \zA{有人说他通中国的诗书，会讲 ‘五经’，能做诗填词。}
}
{
    \eA{In a belt, she carried a Japanese sword, also inlaid with gold and studded
        with precious gems. In very truth, even in pictures, there is no one as
        beautiful as she.}
}
{
}

\Verse{124}
{
    \zA{因此我父亲央烦了一位通官，烦他写了一张字，就写他做 的诗。”}
}
{
    \eA{Some people said that she was thoroughly conversant with Chinese literature,
        and could explain the 'Five classics,' that she was able to write odes and
        devise roundelays, and so my father requested an interpreter to ask her to
        write something.}
}
{
}

\Verse{125}
{
    \zA{众人都称道奇异。}
}
{
    \eA{She thereupon wrote an original stanza, which all, with one voice,}
}
{
}

\Verse{126}
{
    \zA{宝玉忙笑道：“好妹妹，你拿出来我们瞧瞧。”}
}
{
    \eA{praised for its remarkable beauty, and extolled for its extraordinary
        merits." "My dear cousin," eagerly smiled Pao-yü,}
}
{
}

\Verse{127}
{
    \zA{宝琴笑 道：“在南京收着呢，此时那里去取？”}
}
{
    \eA{"produce what she wrote, and let's have a look at it." "It's put away in
        Nanking;" Pao-ch'in replied with a smile.}
}
{
}

\Verse{128}
{
    \zA{宝玉听了，大失所望，便说：“没福得见 这世面！”}
}
{
    \eA{"So how could I at present go and fetch it?" Great was Pao-yü's
        disappointment at this rejoinder. "I've no luck," he cried,}
}
{
}

\Verse{129}
{
    \zA{黛玉笑拉宝琴道：“你别哄我们：我知道你这一来，你的这些东西未必 放在家里，自然都是要带上来的。}
}
{
    \eA{"to see anything like this in the world." Tai-yü laughingly laid hold of
        Pao-ch'in. "Don't be humbugging us!" she remarked. "I know well enough that
        you are not likely, on a visit like this, to have left any such things of
        yours at home. You must have brought them along.}
}
{
}

\Verse{130}
{
    \zA{这会子又扯谎，说没带来。}
}
{
    \eA{Yet here you are now again palming off a fib on us by saying that you
        haven't got them with you.}
}
{
}

\Verse{131}
{
    \zA{他们虽信，我是不信 的。”}
}
{
    \eA{You people may believe what she says, but I, for my part, don't."}
}
{
}

\Verse{132}
{
    \zA{宝琴便红了脸，低头微笑不答。}
}
{
    \eA{Pao-ch'in got red in the face. Drooping her head against her chest, she gave
        a faint smile; but she uttered not a word by way of response.}
}
{
}

\Verse{133}
{
    \zA{宝钗笑道：“偏这颦儿惯说这些话，你就伶 俐的太过了。”}
}
{
    \eA{"Really P'in Erh you've got into the habit of talking like this!" Pao-ch'ai
        laughed. "You're too shrewd by far." "Bring them along," Tai-yü urged with a
        smile,}
}
{
}

\Verse{134}
{
    \zA{黛玉笑道：“带了来，就给我们见识见识也罢了。”}
}
{
    \eA{"and give us a chance of seeing something and learning something; it won't
        hurt them." "There's a whole heap of trunks and baskets,"}
}
{
}

\Verse{135}
{
    \zA{宝钗笑道：“箱 子笼子一大堆，还没理清呢，知道在那个里头呢？}
}
{
    \eA{Pao-ch'ai put in laughing, "which haven't been yet cleared away. And how
        could one tell in which particular one, they're packed up? Wait a few days,}
}
{
}

\Verse{136}
{
    \zA{等过日子收拾清了找出来，大家 再看罢了。”}
}
{
    \eA{and when things will have been put straight a bit, we'll try and find them:
        and every one of us can then have a look at them;}
}
{
}

\Verse{137}
{
    \zA{又向宝琴道：“你要记得，何不念念我们听听？”}
}
{
    \eA{that will be all right. But if you happen to remember the lines," she
        pursued, speaking to Pao-ch'in, "why not recite them for our benefit?"}
}
{
}

\Verse{138}
{
    \zA{宝琴答道：“记得 他做的五言律一首，要论外国的女子，也就难为他了。”}
}
{
    \eA{"I remember so far that her lines consisted of a stanza with five characters
        in each line," Pao-ch'ai returned for answer. "For a foreign girl, they're
        verily very well done." "Don't begin for a while,"}
}
{
}

\Verse{139}
{
    \zA{宝钗道：“你且别念，等 我把云儿叫了来，也叫他听听。”}
}
{
    \eA{Pao-ch'ai exclaimed. "Let me send for Yün Erh, so that she too might hear
        them." After this remark, she called Hsiao Lo to her. "Go to my place," she
        observed,}
}
{
}

\Verse{140}
{
    \zA{说着，便叫小螺来，吩咐道：“你到我那里去， 就说我们这里有一个外国的美人来了，做的好诗，请你这‘诗疯子’来瞧去，再把 我们‘诗呆子’也带来。”}
}
{
    \eA{"and tell her that a foreign beauty has come over, who's a splendid hand at
        poetry. 'You, who have poetry on the brain,' (say to her), 'are invited to
        come and see her,' and then lay hold of this verse-maniac of ours and bring
        her along."}
}
{
}

\Verse{141}
{
    \zA{小螺笑着去了。}
}
{
    \eA{Hsiao Lo gave a smile, and went away.}
}
{
}

\Verse{142}
{
    \zA{半日，只听湘云笑问：“那一个外国的美人来了？”}
}
{
    \eA{After a long time, they heard Hsiang-yün laughingly inquire, "What foreign
        beauty has come?"}
}
{
}

\Verse{143}
{
    \zA{一头说，一头走，和香菱 来了。}
}
{
    \eA{But while asking this question, she made her appearance in company with
        Hsiang Ling.}
}
{
}

\Verse{144}
{
    \zA{众人笑道：“人未见形，先已闻声。”}
}
{
    \eA{"We heard your voices long before we caught a glimpse of your persons!"}
}
{
}

\Verse{145}
{
    \zA{宝琴等让坐，遂把方才的话重告诉了 一遍。}
}
{
    \eA{the party laughed. Pao-ch'in and her companions motioned to her to sit down,
        and, in due course, she reiterated what she had told them a short while
        back.}
}
{
}

\Verse{146}
{
    \zA{湘云笑道：“快念来听听。”}
}
{
    \eA{"Be quick, out with it! Let's hear what it is!" Hsiang-yün smilingly cried.}
}
{
}

\Verse{147}
{
    \zA{宝琴因念道： 昨夜朱楼梦，今宵水国吟。}
}
{
    \eA{Pao-ch'in thereupon recited: Last night in the Purple Chamber I dreamt. This
        evening on the 'Shui Kuo' Isle I sing.}
}
{
}

\Verse{148}
{
    \zA{岛云蒸大海，岚气接丛林。}
}
{
    \eA{The clouds by the isle cover the broad sea. The zephyr from the peaks
        reaches the woods.}
}
{
}

\Verse{149}
{
    \zA{月本无今古，情缘自浅深。}
}
{
    \eA{The moon has never known present or past. From shallow and deep causes
        springs love's fate.}
}
{
}

\Verse{150}
{
    \zA{汉南春历历，焉得不关心？}
}
{
    \eA{When I recall my springs south of the Han, Can I not feel disconsolate at
        heart?}
}
{
}

\Verse{151}
{
    \zA{众人听了，都道：“难为他!竟比我们中国人还强。”}
}
{
    \eA{After listening to her, "She does deserve credit," they unanimously shouted,
        "for she really is far superior to us, Chinese though we be." But scarcely
        was this remark out of their lips,}
}
{
}

\Verse{152}
{
    \zA{一语未了，只见麝月走 来，说：“太太打发了人来告诉二爷，明儿一早往舅舅那里去，就说太太身上不大 好，不得亲身来。”}
}
{
    \eA{when they perceived She Yüeh walk in. "Madame Wang," she said, "has sent a
        servant to inform you, Master Secundus, that 'you are to go at an early hour
        to-morrow morning to your maternal uncle's, and that you are to explain to
        him that her ladyship isn't feeling quite up to the mark, and that she
        cannot pay him a visit in person.'"}
}
{
}

\Verse{153}
{
    \zA{宝玉忙站起来答应道：“是。”}
}
{
    \eA{Pao-yü precipitately jumped to his feet (out of deference to his mother),}
}
{
}

\Verse{154}
{
    \zA{因问宝钗宝琴：“你们二位可 去？”}
}
{
    \eA{and signified his assent, by answering 'Yes.' He then went on to inquire of
        Pao-ch'ai and Pao-ch'in,}
}
{
}

\Verse{155}
{
    \zA{宝钗道：“我们不去。}
}
{
    \eA{"Are you two going?" "We're not going,"}
}
{
}

\Verse{156}
{
    \zA{昨儿单送了礼去了。”}
}
{
    \eA{Pao-ch'ai rejoined. "We simply went there yesterday to take our presents
        over but we left after a short chat."}
}
{
}

\Verse{157}
{
    \zA{大家说了一回方散。}
}
{
    \eA{Pao-yü thereupon pressed his female cousins to go ahead and he then followed
        them.}
}
{
}

\Verse{158}
{
    \zA{宝玉因让诸姐妹先行，自己在后面。}
}
{
    \eA{But Tai-yü called out to him again and stopped him.}
}
{
}

\Verse{159}
{
    \zA{黛玉便又叫住他，问道：“袭人到底多早 晚回来？”}
}
{
    \eA{"When is Hsi Jen, after all, coming back?" she asked. "She'll naturally come
        back after she has accompanied the funeral,"}
}
{
}

\Verse{160}
{
    \zA{宝玉道：“自然等送了殡才来呢。”}
}
{
    \eA{Pao-yü retorted. Tai-yü had something more she would have liked to tell him,
        but she found it difficult to shape it into words.}
}
{
}

\Verse{161}
{
    \zA{黛玉还有话说，又不能出口，出了 一回神，便说道：“你去罢。”}
}
{
    \eA{After some moments spent in abstraction, "Off with you!" she cried. Pao-yü
        too felt that he treasured in his heart many things he would fain confide to
        her,}
}
{
}

\Verse{162}
{
    \zA{宝玉也觉心里有许多话，只是口里不知要说什么， 想了一想，也笑道：“明儿再说罢。”}
}
{
    \eA{but he did not know what to bring to his lips, so after cogitating within
        himself for a time, he likewise observed smilingly: "We'll have another chat
        to-morrow," and, as he said so, he wended his way down the stairs.}
}
{
}

\Verse{163}
{
    \zA{一面下台阶，低头正欲迈步，复又忙回身问 道：“如今夜越发长了，你一夜咳嗽几次？}
}
{
    \eA{Lowering his head, he was just about to take a step forward, when he twisted
        himself round again with alacrity. "Now that the nights are longer than they
        were, you're sure to cough often and wake several times in the night;}
}
{
}

\Verse{164}
{
    \zA{醒几遍？”}
}
{
    \eA{eh?" he asked. "Last night,"}
}
{
}

\Verse{165}
{
    \zA{黛玉道：“昨儿夜里好了， 只咳嗽两遍，却只睡了四更一个更次，就再不能睡了。”}
}
{
    \eA{Tai-yü answered, "I was all right; I coughed only twice. But I only slept at
        the fourth watch for a couple of hours and then I couldn't close my eyes
        again." "I really have something very important to tell you,"}
}
{
}

\Verse{166}
{
    \zA{宝玉又笑道：“正是有句 要紧的话，这会子才想起来。”}
}
{
    \eA{Pao-yü proceeded with another smile. "It only now crossed my mind." Saying
        this, he approached her and added in a confidential tone:}
}
{
}

\Verse{167}
{
    \zA{一面说，一面便挨近身来，悄悄道：“我想宝姐姐 送你的燕窝――”一语未了，只见赵姨娘走进来瞧黛玉，问：“姑娘这几天可好了？”}
}
{
    \eA{"I think that the birds' nests sent to you by cousin Pao-chai…." Barely,
        however, had he had time to conclude than he spied dame Chao enter the room
        to pay Tai-yü a visit. "Miss, have you been all right these last few days?"
        she inquired. Tai-yü readily guessed that this was an attention extended to
        her merely as she had,}
}
{
}

\Verse{168}
{
    \zA{黛玉便知他从探春处来，从门前过，顺路的人情，忙陪笑让坐，说：“难得姨娘想 着，怪冷的，亲自走来。”}
}
{
    \eA{on her way back from T'an Ch'un's quarters, to pass by her door, so speedily
        smiling a forced smile, she offered her a seat. "Many thanks, dame Chao,"
        she said, "for the trouble of thinking of me, and for coming in person in
        this intense cold."}
}
{
}

\Verse{169}
{
    \zA{又忙命倒茶，一面又使眼色给宝玉。}
}
{
    \eA{Hastily also bidding a servant pour the tea, she simultaneously winked at
        Pao-yü.}
}
{
}

\Verse{170}
{
    \zA{宝玉会意，便走了 出来。}
}
{
    \eA{Pao-yü grasped her meaning, and forthwith quitted the apartment.}
}
{
}

\Verse{171}
{
    \zA{正值吃晚饭时，见了王夫人，又嘱咐他早去。}
}
{
    \eA{As this happened to be about dinner time, and he had been enjoined as well
        by Madame Wang to be back at an early hour,}
}
{
}

\Verse{172}
{
    \zA{宝玉回来，看晴雯吃了药。}
}
{
    \eA{Pao-yü returned to his quarters, and looked on while Ch'ing Wen took her
        medicine.}
}
{
}

\Verse{173}
{
    \zA{此 夕宝玉便不命晴雯挪出暖阁来，自己便在晴雯外边。}
}
{
    \eA{Pao-yü did not desire Ch'ing Wen this evening to move into the winter
        apartment, but stayed with Ch'ing Wen outside; and, giving orders to bring
        the warming-frame near the winter apartment,}
}
{
}

\Verse{174}
{
    \zA{又命将熏笼抬至暖阁前，麝月 便在熏笼上睡。}
}
{
    \eA{She Yueh slept on it. Nothing of any interest worth putting on record
        transpired during the night. On the morrow, before the break of day, Ch'ing
        Wen aroused She Yueh.}
}
{
}

\Verse{175}
{
    \zA{一宿无话。}
}
{
    \eA{"You should awake," she said.}
}
{
}

\Verse{176}
{
    \zA{至次日天未明，晴雯便叫醒麝月道：“你也该醒了，只是睡不够。}
}
{
    \eA{"The only thing is that you haven't had enough sleep. If you go out and tell
        them to get the water for tea ready for him, while I wake him, it will be
        all right." She Yueh immediately jumped up and threw something over her.}
}
{
}

\Verse{177}
{
    \zA{你出去叫人 给他预备茶水，我叫醒他就是了。”}
}
{
    \eA{"Let's call him to get up and dress in his fine clothes." she said. "We can
        summon them in, after this fire-box has been removed.}
}
{
}

\Verse{178}
{
    \zA{麝月忙披衣起来道：“咱们叫他起来，穿好衣 裳，抬过这火箱去，再叫他们进来。}
}
{
    \eA{The old nurses told us not to allow him to stay in this room for fear the
        virus of the disease should pass on to him; so now if they see us bundled up
        together in one place, they're bound to kick up another row." "That's my
        idea too," Ch'ing Wen replied.}
}
{
}

\Verse{179}
{
    \zA{老妈妈们已经说过，不叫他在这屋里，怕过了 病气；如今他们见咱们挤在一处，又该唠叨了。”}
}
{
    \eA{The two girls were then about to call him, when Pao-yü woke up of his own
        accord, and speedily leaping out of bed, he threw his clothes over him. She
        Yüeh first called a young maid into the room and put things shipshape before
        she told Ch'in Wen and the other servant-girls to enter;}
}
{
}

\Verse{180}
{
    \zA{晴雯道：“我也是这么说。”}
}
{
    \eA{and along with them, she remained in waiting upon Pao-yü while he combed his
        hair, and washed his face and hands.}
}
{
}

\Verse{181}
{
    \zA{二 人才叫时，宝玉已醒了，忙起身披衣。}
}
{
    \eA{This part of his toilet over, She Yüeh remarked: "It's cloudy again, so I
        suppose it's going to snow.}
}
{
}

\Verse{182}
{
    \zA{麝月先叫进小丫头子来收拾妥了，才命秋纹 等进来，一同伏侍。}
}
{
    \eA{You'd better therefore wear a woollen overcoat!" Pao-yü nodded his head
        approvingly; and set to work at once to effect the necessary change in his
        costume. A young waiting-maid then presented him a covered bowl,}
}
{
}

\Verse{183}
{
    \zA{宝玉梳洗已毕，麝月道：“天又阴阴的，只怕下雪，穿一套毡 子的罢。”}
}
{
    \eA{in a small tea tray, containing a decoction made of Fu-kien lotus and red
        dates. After Pao-yü had had a couple of mouthfuls, She Yüeh also brought him
        a small plateful of brown ginger,}
}
{
}

\Verse{184}
{
    \zA{宝玉点头，即时换了衣裳。}
}
{
    \eA{prepared according to some prescription. Pao-yü put a piece into his mouth,}
}
{
}

\Verse{185}
{
    \zA{小丫头便用小茶盘捧了一盖碗建莲红枣汤来， 宝玉喝了两口；麝月又捧过一小碟法制紫姜来，宝玉噙了一块。}
}
{
    \eA{and, impressing some advice on Ch'ing 'Wen, he crossed over to dowager lady
        Chia's suite of rooms. His grandmother had not yet got out of bed. But she
        was well aware that Pao-yü was going out of doors so having the entrance
        leading into her bedroom opened she asked Pao-yü to walk in.}
}
{
}

\Verse{186}
{
    \zA{又嘱咐了晴雯，便 忙往贾母处来。}
}
{
    \eA{Pao-yü espied behind the old lady, Pao-ch'in lying with her face turned
        towards the inside,}
}
{
}

\Verse{187}
{
    \zA{贾母犹未起来，知道宝玉出门，便开了屋门，命宝玉进去。}
}
{
    \eA{and not awake yet from her sleep. Dowager lady Chia observed that Pao-yü was
        clad in a deep-red felt fringed overcoat, with woollen lichee-coloured
        archery-sleeves and with an edging of dark green glossy satin,}
}
{
}

\Verse{188}
{
    \zA{宝玉见贾母身后宝 琴面向里睡着未醒。}
}
{
    \eA{embroidered with gold rings. "What!" old lady Chia inquired, "is it
        snowing?" "The weather is dull,"}
}
{
}

\Verse{189}
{
    \zA{贾母见宝玉身上穿着荔支色哆罗呢的箭袖，大红猩猩毡盘金彩 绣石青妆缎沿边的排穗褂。}
}
{
    \eA{Pao-yü replied, "but it isn't snowing yet." Dowager lady Chia thereupon sent
        for Yüan Yang and told her to fetch the peacock down pelisse, finished the
        day before, and give it to him. Yüan Yang signified her obedience and went
        off,}
}
{
}

\Verse{190}
{
    \zA{贾母道：“下雪呢么？”}
}
{
    \eA{and actually returned with what was wanted. When Pao-yü came to survey it,}
}
{
}

\Verse{191}
{
    \zA{宝玉道：“天阴着，还没下呢。”}
}
{
    \eA{he found that the green and golden hues glistened with bright lustre,}
}
{
}

\Verse{192}
{
    \zA{贾母便命：“鸳鸯来，把昨儿那一件孔雀毛的氅衣给他罢。”}
}
{
    \eA{that the jadelike variegated colours on it shone with splendour, and that it
        bore no resemblance to the duck-down coat, which Pao-ch'in had been wearing.
        "This," he heard his grandmother smilingly remark,}
}
{
}

\Verse{193}
{
    \zA{鸳鸯答应走去，果取 了一件来。}
}
{
    \eA{"is called 'bird gold'. This is woven of the down of peacocks, caught in
        Russia, twisted into thread.}
}
{
}

\Verse{194}
{
    \zA{宝玉看时，金翠辉煌，碧彩？}
}
{
    \eA{The other day, I presented that one with the wild duck down to your young
        female cousin, so I now give you this one."}
}
{
}

\Verse{195}
{
    \zA{灼，又不似宝琴所披之凫靥裘。}
}
{
    \eA{Pao-yü prostrated himself before her, after which he threw the coat over his
        shoulders.}
}
{
}

\Verse{196}
{
    \zA{只听贾母 笑道：“这叫做‘雀金呢’，这是俄罗斯国拿孔雀毛拈了线织的。}
}
{
    \eA{"Go and let your mother see it before you start," his grandmother laughingly
        added. Pao-yü assented, and quitted her apartments, when he caught sight of
        Yüan Yang standing below rubbing her eyes.}
}
{
}

\Verse{197}
{
    \zA{前儿那件野鸭子 的给了你小妹妹，这件给你罢。”}
}
{
    \eA{Ever since the day on which Yüan Yang had sworn to have done with the match,
        she had not exchanged a single word with Pao-yü. Pao-yü was therefore day
        and night a prey to dejection.}
}
{
}

\Verse{198}
{
    \zA{宝玉磕了一个头，便披在身上。}
}
{
    \eA{So when he now observed her shirk his presence again, Pao-yü at once
        advanced up to her, and, putting on a smile,}
}
{
}

\Verse{199}
{
    \zA{贾母笑道：“你 先给你娘瞧瞧去再去。”}
}
{
    \eA{"My dear girl," he said, "do look at the coat I've got on. Is it nice or
        not?" Yüan Yang shoved his hand away,}
}
{
}

\Verse{200}
{
    \zA{宝玉答应了，便出来，只见鸳鸯站在地下揉眼睛。}
}
{
    \eA{and promptly walked into dowager lady Chia's quarters. Pao-yü was thus
        compelled to repair to Madame Wang's room, and let her see his coat.}
}
{
}

\Verse{201}
{
    \zA{因自那 日鸳鸯发誓绝婚之后，他总不合宝玉说话，宝玉正自日夜不安，此时见他又要回避， 宝玉便上来笑道：“好姐姐你瞧瞧，我穿着这个好不好？”}
}
{
    \eA{Retracing afterwards his footsteps into the garden, he let Ch'ing Wen and
        She Yüeh also have a look at it, and then came and told his grandmother that
        he had attended to her wishes. "My mother," he added, "has seen what I've
        got on. But all she said was: 'what a pity!' and then she went on to enjoin
        me to be 'careful with it and not to spoil it.'"}
}
{
}

\Verse{202}
{
    \zA{鸳鸯一摔手，便进贾母 屋里来了。}
}
{
    \eA{"There only remains this single one," old lady Chia observed, "so if you
        spoil it you can't have another.}
}
{
}

\Verse{203}
{
    \zA{宝玉只得到了王夫人屋里，给王夫人看了，然后又回至园中，给晴雯麝 月看过，来回覆贾母说：“太太看了，只说可惜了的，叫我仔细穿，别遭塌了。”}
}
{
    \eA{Even did I want to have one made for you like it now, it would be out of the
        question." At the close of these words, she went on to advise him. "Don't,"
        she said, "have too much wine and come back early." Pao-yü acquiesced by
        uttering several yes's.}
}
{
}

\Verse{204}
{
    \zA{贾母道：“就剩了这一件，你遭塌了也再没了。}
}
{
    \eA{An old nurse then followed him out into the pavilion. Here they discovered
        six attendants, (that is), Pao-yü's milk-brother Li Kuei,}
}
{
}

\Verse{205}
{
    \zA{这会子特给你做这个，也是没有的 事。”}
}
{
    \eA{and Wang Ho-jung, Chang Jo-chin, Chao I-hua, Ch'ien Ch'i, and Chou Jui, as
        well as four young servant-lads:}
}
{
}

\Verse{206}
{
    \zA{说着又嘱咐：“不许多吃酒，早些回来。”}
}
{
    \eA{Pei Ming, Pan Ho, Chu Shao and Sao Hung; some carrying bundles of clothes on
        their backs, some holding cushions in their hands,}
}
{
}

\Verse{207}
{
    \zA{宝玉应了几个“是”。}
}
{
    \eA{others leading a white horse with engraved saddle and variegated bridles.}
}
{
}

\Verse{208}
{
    \zA{老嬷嬷跟至厅上，只见宝玉的奶兄李贵、王荣和张若锦、 赵亦华、钱升、周瑞六个人，带着焙茗、伴鹤、锄药、扫红四个小厮，背着衣包，
        拿着坐褥，笼着一匹雕鞍彩辔的白马，已伺候多时了。}
}
{
    \eA{They had already been waiting for a good long while. The old nurse went on
        to issue some directions, and the six servants, hastily expressing their
        obedience by numerous yes's, quickly caught hold of the saddle and weighed
        the stirrup down while Pao-yü mounted leisurely. Li Kuei and Wang Ho-jung
        then led the horse by the bit. Two of them, Ch'ien Ch'i and Chou Jui,}
}
{
}

\Verse{209}
{
    \zA{老嬷嬷又嘱咐他们些话，六 个人连应了几个“是”，忙捧鞍坠镫，宝玉慢慢的上了马。}
}
{
    \eA{walked ahead and showed the way. Chang Jo-chin and Chao I-hua followed
        Pao-yü closely on each side. "Brother Chou and brother Ch'ien," Pao-yü
        smiled, from his seat on his horse, "let's go by this side-gate.}
}
{
}

\Verse{210}
{
    \zA{李贵王荣笼着嚼环，钱 升周瑞二人在前引导，张若锦赵亦华在两边，紧贴宝玉身后。}
}
{
    \eA{It will save my having again to dismount, when we reach the entrance to my
        father's study." "Mr. Chia Cheng is not in his study," Chou Jui laughed,
        with a curtsey. "It has been daily under lock and key, so there will be no
        need for you,}
}
{
}

\Verse{211}
{
    \zA{宝玉在马上笑道：“周 哥，钱哥，咱们打这角门走罢，省了到老爷的书房门口，又下来。”}
}
{
    \eA{master, to get down from your horse." "Though it be locked up," Pao-yü
        smiled, "I shall have to dismount all the same." "You're quite right in what
        you say, master;" both Ch'ien Ch'i and Li Kuei chimed in laughingly;}
}
{
}

\Verse{212}
{
    \zA{周瑞侧身笑道： “老爷不在书房里，天天锁着，爷可以不用下来罢了。”}
}
{
    \eA{"but pretend you're lazy and don't get down. In the event of our coming
        across Mr. Lai Ta and our number two Mr. Lin, they're sure, rather awkward
        though it be for them to say anything to their master,}
}
{
}

\Verse{213}
{
    \zA{宝玉笑道：“虽锁着，也 要下来的。”}
}
{
    \eA{to tender you one or two words of advice, but throw the whole of the blame
        upon us.}
}
{
}

\Verse{214}
{
    \zA{钱升李贵都笑道：“爷说的是。}
}
{
    \eA{You can also tell them that we had not explained to you what was the right
        thing to do."}
}
{
}

\Verse{215}
{
    \zA{就托懒不下来，倘或遇见赖大爷林二 爷，虽不好说爷，也要劝两句。}
}
{
    \eA{Chou Jui and Ch'ien Ch'i accordingly wended their steps straight for the
        side-gate. But while they were keeping up some sort of conversation, they
        came face to face with Lai Ta on his way in.}
}
{
}

\Verse{216}
{
    \zA{所有的不是，都派在我们身上，又说我们不教给爷 礼了。”}
}
{
    \eA{Pao-yü speedily pulled in his horse, with the idea of dismounting. But Lai
        Ta hastened to draw near and to clasp his leg. Pao-yü stood up on his
        stirrup,}
}
{
}

\Verse{217}
{
    \zA{周瑞钱升便一直出角门来。}
}
{
    \eA{and, putting on a smile, he took his hand in his, and made several remarks
        to him.}
}
{
}

\Verse{218}
{
    \zA{正说话时，顶头见赖大进来，宝玉忙笼住马，意欲下来。}
}
{
    \eA{In quick succession, he also perceived a young servant-lad make his
        appearance inside leading the way for twenty or thirty servants, laden with
        brooms and dust-baskets.}
}
{
}

\Verse{219}
{
    \zA{赖大忙上来抱住腿。}
}
{
    \eA{The moment they espied Pao-yü, they, one and all, stood along the wall,}
}
{
}

\Verse{220}
{
    \zA{宝玉便在镫上站起来，笑着，携手说了几句话。}
}
{
    \eA{and dropped their arms against their sides, with the exception of the head
        lad, who bending one knee, said: "My obeisance to you, sir."}
}
{
}

\Verse{221}
{
    \zA{接着又见个小厮带着二三十人，拿 着扫帚簸箕进来，见了宝玉，都顺墙垂手立住，独为首的小厮打了个千儿，说：“请 爷安。”}
}
{
    \eA{Pao-yü could not recall to mind his name or surname, but forcing a faint
        smile, he nodded his head to and fro. It was only when the horse had well
        gone past, that the lad eventually led the bevy of servants off, and that
        they went after their business. Presently, they egressed from the side-gate.}
}
{
}

\Verse{222}
{
    \zA{宝玉不知名姓，只微笑点点头儿。}
}
{
    \eA{Outside, stood the servant-lads of the six domestics, Li Kuei and his
        companions, as well as several grooms,}
}
{
}

\Verse{223}
{
    \zA{马已过去，那人方带人去了。}
}
{
    \eA{who had, from an early hour, got ready about ten horses and been standing,}
}
{
}

\Verse{224}
{
    \zA{于是出了 角门。}
}
{
    \eA{on special duty, waiting for their arrival.}
}
{
}

\Verse{225}
{
    \zA{外有李贵等六人的小厮并几个马夫，早预备下十来匹马专候，一出角门，李 贵等各上马前引，一阵烟去了，不在话下。}
}
{
    \eA{As soon as they reached the further end of the side-gate, Li Kuei and each
        of the other attendants mounted their horses, and pressed ahead to lead the
        way. Like a streak of smoke, they got out of sight, without any occurrence
        worth noticing. Ch'ing Wen, meanwhile, continued to take her medicines.}
}
{
}

\Verse{226}
{
    \zA{这里晴雯吃了药仍不见病退，急的乱骂大夫，说：“只会哄人的钱，一剂好药 也不给人吃。”}
}
{
    \eA{But still she experienced no relief in her ailment. Such was the state of
        exasperation into which she worked herself that she abused the doctor right
        and left. "All he's good for," she cried, "is to squeeze people's money. But
        he doesn't know how to prescribe a single dose of efficacious medicine for
        his patients."}
}
{
}

\Verse{227}
{
    \zA{麝月笑劝他道：“你太性急了，俗语说：‘病来如山倒，病去如抽 丝。’}
}
{
    \eA{"You have far too impatient a disposition!" She Yüeh said, as she advised
        her, with a smile. "'A disease,' the proverb has it, 'comes like a crumbling
        mountain, and goes like silk that is reeled.'}
}
{
}

\Verse{228}
{
    \zA{又不是老君的仙丹，那有这么灵药？}
}
{
    \eA{Besides, they're not the divine pills of 'Lao Chün'. How ever could there be
        such efficacious medicines?}
}
{
}

\Verse{229}
{
    \zA{你只静养几天，自然就好了。}
}
{
    \eA{The only thing for you to do is to quietly look after yourself for several
        days, and you're sure to get all right.}
}
{
}

\Verse{230}
{
    \zA{你越急越 着手。”}
}
{
    \eA{But the more you work yourself into such a frenzy,}
}
{
}

\Verse{231}
{
    \zA{晴雯又骂小丫头子们：“那里攒沙去了!瞅着我病了，都大胆子走了。}
}
{
    \eA{the worse you get!" Ch'ing Weng went on to heap abuse on the head of the
        young-maids. "Where have they gone? Have they bored into the sand?" she
        ejaculated. "They see well enough that I'm ill,}
}
{
}

\Verse{232}
{
    \zA{明 儿我好了，一个个的才揭了你们的皮！”}
}
{
    \eA{so they make bold and runaway. But by and bye when I recover, I shall take
        one by one of you and flay your skin off for you."}
}
{
}

\Verse{233}
{
    \zA{唬的小丫头子定儿忙进来问：“姑娘做什 么？”}
}
{
    \eA{Ting Erh, a young maid, was struck with dismay, and ran up to her with hasty
        step. "Miss," she inquired, "what's up with you?"}
}
{
}

\Verse{234}
{
    \zA{晴雯道：“别人都死了，就剩了你不成？”}
}
{
    \eA{"Is it likely that the rest are all dead and gone, and that there only
        remains but you?" Ch'ing Wen exclaimed. But while she spoke,}
}
{
}

\Verse{235}
{
    \zA{说着，只见坠儿也蹭进来了。}
}
{
    \eA{she saw Chui Erh also slowly enter the room. "Look at this vixen!" Ch'ing
        Wen shouted.}
}
{
}

\Verse{236}
{
    \zA{晴 雯道：“你瞧瞧这小蹄子，不问他还不来呢。}
}
{
    \eA{"If I don't ask for her, she won't come. Had there been any monthly
        allowances issued and fruits distributed here, you would have been the first
        to run in!}
}
{
}

\Verse{237}
{
    \zA{这里又放月钱了，又散果子了，你该 跑在头里了。}
}
{
    \eA{But approach a bit! Am I tigress to gobble you up?" Chui Erh was under the
        necessity of advancing a few steps nearer to her. But, all of a sudden,
        Ch'ing Wen stooped forward,}
}
{
}

\Verse{238}
{
    \zA{你往前些!我是老虎，吃了你？”}
}
{
    \eA{and with a dash clutching her hand, she took a long pin from the side of her
        pillow,}
}
{
}

\Verse{239}
{
    \zA{坠儿只得往前凑了几步。}
}
{
    \eA{and pricked it at random all over. "What's the use of such paws?"}
}
{
}

\Verse{240}
{
    \zA{晴雯便冷 不防欠身，一把将他的手抓住，向枕边拿起一丈青来，向他手上乱戳，又骂道：“要 这爪子做什么？}
}
{
    \eA{she railed at her. "They don't ply a needle, and they don't touch any
        thread! All you're good for is to prig things to stuff that mouth of yours
        with! The skin of your phiz is shallow and those paws of yours are light!
        But with the shame you bring upon yourself before the world,}
}
{
}

\Verse{241}
{
    \zA{拈不动针，拿不动线，只会偷嘴吃!眼皮子又浅，爪子又轻，打嘴现 世的，不如戳烂了！”}
}
{
    \eA{isn't it right that I should prick you, and make mincemeat of you?" Chui Erh
        shouted so wildly from pain that She Yueh stepped forward and immediately
        drew them apart. She then pressed Ch'ing Wen,}
}
{
}

\Verse{242}
{
    \zA{坠儿疼的乱喊。}
}
{
    \eA{until she induced her to lie down. "You're just perspiring,"}
}
{
}

\Verse{243}
{
    \zA{麝月忙拉开，按着晴雯躺下，道：“你才出 了汗，又作死!等你好了，要打多少打不得？}
}
{
    \eA{she remarked, "and here you are once more bent upon killing yourself. Wait
        until you are yourself again! Won't you then be able to give her as many
        blows as you may like? What's the use of kicking up all this fuss just now?"}
}
{
}

\Verse{244}
{
    \zA{这会子闹什么。”}
}
{
    \eA{Ch'ing Wen bade a servant tell nurse Sung to come in.}
}
{
}

\Verse{245}
{
    \zA{晴雯便命人叫宋嬷嬷进来，说道：“宝二爷才告诉了我，叫我告诉你们，坠儿 很懒，宝二爷当面使他，他拨嘴儿不动，连袭人使他，他也背地里骂。}
}
{
    \eA{"Our master Secundus, Mr. Pao-yü, recently asked me to tell you," she
        remarked on her arrival, "that Chui Erh is very lazy. He himself gives her
        orders to her very face, but she is ever ready to raise objections and not
        to budge. Even when Hsi Jen bids her do things,}
}
{
}

\Verse{246}
{
    \zA{今儿务必打 发他出去，明儿宝二爷亲自回太太就是了。”}
}
{
    \eA{she vilifies her behind her back. She must absolutely therefore be packed
        off to-day. And if Mr. Pao himself lays the matter to-morrow before Madame
        Wang,}
}
{
}

\Verse{247}
{
    \zA{宋嬷嬷听了，心下便知镯子事发，因 笑道：“虽如此说，也等花姑娘回来，知道了，再打发他。”}
}
{
    \eA{things will be square." After listening to her grievances, nurse Sung
        readily concluded in her mind that the affair of the bracelet had come to be
        known. "What you suggest is well and good, it's true," she consequently
        smiled,}
}
{
}

\Verse{248}
{
    \zA{晴雯道：“宝二爷今 儿千叮咛万嘱咐的，什么‘花姑娘’‘草姑娘’的，我们自然有道理!你只依我的 话，快叫他家的人来领他出去。”}
}
{
    \eA{"but it's as well to wait until Miss Hua (flower) returns and hears about
        the things. We can then give her the sack." "Mr. Pao-yü urgently enjoined
        this to-day," Ch'ing Wen pursued, "so what about Miss Hua (flower) and Miss
        Ts'ao (grass)?}
}
{
}

\Verse{249}
{
    \zA{麝月道：“这也罢了。}
}
{
    \eA{We've, of course, gob rules of propriety here,}
}
{
}

\Verse{250}
{
    \zA{早也是去，晚也是去，早 带了去，早清净一日。”}
}
{
    \eA{so you just do as I tell you; and be quick and send for some one from her
        house to come and fetch her away!" "Well, now let's drop this!"}
}
{
}

\Verse{251}
{
    \zA{宋嬷嬷听了，只得出去唤了他母亲来，打点了他的东西。}
}
{
    \eA{She Yüeh interposed. "Whether she goes soon or whether she goes late is one
        and the same thing; so let them take her away soon; we'll then be the sooner
        clear of her."}
}
{
}

\Verse{252}
{
    \zA{又见了晴雯等，说道：“姑娘们怎么了？}
}
{
    \eA{At these words, nurse Sung had no alternative but to step out, and to send
        for her mother. When she came, she got ready all her effects,}
}
{
}

\Verse{253}
{
    \zA{你侄女儿不好，你们教导他，怎么撵出去？}
}
{
    \eA{and then came to see Ch'ing Wen and the other girls. "Young ladies," she
        said, "what's up? If your niece doesn't behave as she ought to,}
}
{
}

\Verse{254}
{
    \zA{也到底给我们留个脸儿。”}
}
{
    \eA{why, call her to account. But why banish her from this place?}
}
{
}

\Verse{255}
{
    \zA{晴雯道：“这话只等宝玉来问他，与我们无干。”}
}
{
    \eA{You should, indeed, leave us a little face!" "As regards what you say,"
        Ch'ing Wen put in, "wait until Pao-yü comes, and then we can ask him.}
}
{
}

\Verse{256}
{
    \zA{那媳 妇冷笑道：“我有胆子问他去？}
}
{
    \eA{It's nothing to do with us." The woman gave a sardonic smile. "Have I got
        the courage to ask him?"}
}
{
}

\Verse{257}
{
    \zA{他那一件事不是听姑娘们的调停？}
}
{
    \eA{she answered. "In what matter doesn't he lend an ear to any settlement you,
        young ladies, may propose?}
}
{
}

\Verse{258}
{
    \zA{他纵依了，姑娘们 不依，也未必中用。}
}
{
    \eA{He invariably agrees to all you say! But if you, young ladies, aren't
        agreeable, it's really of no avail.}
}
{
}

\Verse{259}
{
    \zA{比如方才说话，虽背地里，姑娘就直叫他的名字，在姑娘们就 使得，在我们就成了野人了！”}
}
{
    \eA{When you, for example, spoke just now,—it's true it was on the sly,—you
        called him straightway by his name, miss. This thing does very well with
        you, young ladies, but were we to do anything of the kind, we'd be looked
        upon as very savages!"}
}
{
}

\Verse{260}
{
    \zA{晴雯听说，越发急红了脸，说道：“我叫了他的名字了。}
}
{
    \eA{Ch'ing Wen, upon hearing her remark, became more than ever exasperated, and
        got crimson in the face. "Yes, I called him by his name,"}
}
{
}

\Verse{261}
{
    \zA{你在老太太、太太跟 前告我去，说我野，也撵出我去！”}
}
{
    \eA{she rejoined, "so you'd better go and report me to our old lady and Madame
        Wang. Tell them I'm a rustic and let them send me too off." "Sister-in-law,"
        urged She Yüeh, "just you take her away;}
}
{
}

\Verse{262}
{
    \zA{麝月道：“嫂子你只管带了人出去，有话再说。}
}
{
    \eA{and if you've got aught to say, you can say it by and bye. Is this a place
        for you to bawl in and to try and explain what is right?}
}
{
}

\Verse{263}
{
    \zA{这个地方岂有你叫喊讲理的？}
}
{
    \eA{Whom have you seen discourse upon the rules of propriety with us? Not to
        speak of you, sister-in-law,}
}
{
}

\Verse{264}
{
    \zA{你见谁和我们讲过理？}
}
{
    \eA{even Mrs. Lai Ta and Mrs. Lin treat us fairly well. And as for calling him
        by name,}
}
{
}

\Verse{265}
{
    \zA{别说嫂子你，就是赖大奶奶、林 大娘也得担待我们三分。}
}
{
    \eA{why, from days of yore to the very present, our dowager mistress has
        invariably bidden us do so. You yourselves are well aware of it.}
}
{
}

\Verse{266}
{
    \zA{就是叫名字，从小儿直到如今，都是老太太吩咐过的，你 们也知道的：恐怕难养活，巴巴的写了他的小名儿各处贴着，叫万人叫去，为的是
        好养活，连挑水挑粪花子都叫得，何况我们!连昨儿林大娘叫了一声‘爷’，老太 太还说呢。}
}
{
    \eA{So much did she fear that it would be a difficult job to rear him that she
        deliberately wrote his infant name on slips of paper and had them stuck
        everywhere and anywhere with the design that one and all should call him by
        it. And this in order that it might exercise a good influence upon his
        bringing up. Even water-coolies and scavenger-coolies indiscriminately
        address him by his name; and how much more such as we? So late, in fact, as
        yesterday Mrs. Lin gave him but once the title of 'Sir,' and our old
        mistress called even her to task.}
}
{
}

\Verse{267}
{
    \zA{此是一件。}
}
{
    \eA{This is one side of the question.}
}
{
}

\Verse{268}
{
    \zA{二则我们这些人，常回老太太、太太的话去，可不叫着名回 话，难道也称‘爷’？}
}
{
    \eA{In the next place, we all have to go and make frequent reports to our
        venerable dowager lady and Madame Wang, and don't we with them allude to him
        by name in what we have to say? Is it likely we'd also style him 'Sir?'}
}
{
}

\Verse{269}
{
    \zA{那一日不把‘宝玉’两字叫二百遍，偏嫂子又来挑这个了!过 一天嫂子闲了，在老太太、太太跟前听听我们当着面儿叫他，就知道了。}
}
{
    \eA{What day is there that we don't utter the two words 'Pao-yü' two hundred
        times? And is it for you, sister-in-law, to come and pick out this fault?
        But in a day or so, when you've leisure to go to our old mistress' and
        Madame Wang's, you'll hear us call him by name in their very presence, and
        then you'll feel convinced.}
}
{
}

\Verse{270}
{
    \zA{嫂子原也 不得在老太太、太太跟前当些体统差使，成年家只在三门外头混，怪不得不知道我 们里头的规矩。}
}
{
    \eA{You've never, sister-in-law, had occasion to fulfil any honourable duties by
        our old lady and our lady. From one year's end to the other, all you do is
        to simply loaf outside the third door. So it's no matter of surprise, if you
        don't happen to know anything of the customs which prevail with us inside.}
}
{
}

\Verse{271}
{
    \zA{这里不是嫂子久站的，再一会，不用我们说话，就有人来问你了。}
}
{
    \eA{But this isn't a place where you, sister-in-law, can linger for long. In
        another moment, there won't be any need for us to say anything; for some one
        will be coming to ask you what you want,}
}
{
}

\Verse{272}
{
    \zA{有什么分证的话，且带了他去，你回了林大娘，叫他来找二爷说话。}
}
{
    \eA{and what excuse will you be able to plead? So take her away and let Mrs. Lin
        know about it; and commission her to come and find our Mr. Secundus and tell
        him all. There are in this establishment over a thousand inmates;}
}
{
}

\Verse{273}
{
    \zA{家里上千的人， 他也跑来，我也跑来，我们认人问姓还认不清呢！”}
}
{
    \eA{one comes and another comes, so that though we know people and inquire their
        names, we can't nevertheless imprint them clearly on our minds." At the
        close of this long rigmarole,}
}
{
}

\Verse{274}
{
    \zA{说着，便叫小丫头子：“拿了 擦地的布来擦地！”}
}
{
    \eA{she at once told a young maid to take the mop and wash the floors. The woman
        listened patiently to her arguments, but she could find no words to say
        anything to her by way of reply.}
}
{
}

\Verse{275}
{
    \zA{那媳妇听了，无言可对，亦不敢久站，赌气带了坠儿就走。}
}
{
    \eA{Nor did she have the audacity to protract her stay. So flying into a huff,
        she took Chui Erh along with her, and there and then made her way out. "Is
        it likely," nurse Sung hastily observed,}
}
{
}

\Verse{276}
{
    \zA{宋 嬷嬷忙道：“怪道你这嫂子不知规矩。}
}
{
    \eA{"that a dame like you doesn't know what manners mean? Your daughter has been
        in these rooms for some time, so she should, when she is about to go,}
}
{
}

\Verse{277}
{
    \zA{你女儿在屋里一场，临去时也给姑娘们磕个 头。}
}
{
    \eA{knock her head before the young ladies. She has no other means of showing
        her gratitude. Not that they care much about such things.}
}
{
}

\Verse{278}
{
    \zA{没有别的谢礼，他们也不希罕，不过磕个头尽心罢咧，怎么说走就走？”}
}
{
    \eA{Yet were she to simply knock her head, she would acquit herself of a duty,
        if nothing more. But how is it that she says I'm going, and off she
        forthwith rushes?" Chui Erh overheard these words,}
}
{
}

\Verse{279}
{
    \zA{坠儿 听了，只得翻身进来，给他两个磕头。}
}
{
    \eA{and felt under the necessity of turning back. Entering therefore the
        apartment, she prostrated herself before the two girls,}
}
{
}

\Verse{280}
{
    \zA{又找秋纹等，他们也并不睬他。}
}
{
    \eA{and then she went in quest of Ch'iu Wen and her companions, but neither did
        they pay any notice whatever to her.}
}
{
}

\Verse{281}
{
    \zA{那媳妇？}
}
{
    \eA{"Hai!" ejaculated the woman,}
}
{
}

\Verse{282}
{
    \zA{声 叹气，口不敢言，抱恨而去。}
}
{
    \eA{and heaving a sigh—for she did not venture to utter a word,—she walked off,}
}
{
}

\Verse{283}
{
    \zA{晴雯方才又闪了风，着了气，反觉更不好了。}
}
{
    \eA{fostering a grudge in her heart. Ch'ing Wen had, while suffering from a
        cold, got into a fit of anger into the bargain,}
}
{
}

\Verse{284}
{
    \zA{翻腾至掌灯，刚安静了些，只见 宝玉回来，进门就？}
}
{
    \eA{so instead of being better, she was worse, and she tossed and rolled until
        the time came for lighting the lamps. But the moment she felt more at ease,}
}
{
}

\Verse{285}
{
    \zA{声顿脚。}
}
{
    \eA{she saw Pao-yü come back.}
}
{
}

\Verse{286}
{
    \zA{麝月忙问原故，宝玉道：“今儿老太太喜喜欢欢的给 了这件褂子，谁知不防，后襟子上烧了一块。}
}
{
    \eA{As soon as he put his foot inside the door, he gave way to an exclamation,
        and stamped his foot. "What's the reason of such behaviour?" She Yüeh
        promptly asked him. "My old grandmother," Pao-yü explained, "was in such
        capital spirits that she gave me this coat to-day;}
}
{
}

\Verse{287}
{
    \zA{幸而天晚了，老太太、太太都不理论。”}
}
{
    \eA{but, who'd have thought it, I inadvertently burnt part of the back lapel.
        Fortunately however the evening was advanced so that neither she nor my
        mother noticed what had happened."}
}
{
}

\Verse{288}
{
    \zA{一面脱下来。}
}
{
    \eA{Speaking the while, he took it off.}
}
{
}

\Verse{289}
{
    \zA{麝月瞧时，果然有指顶大的烧眼，说：“这必定是手炉里的火迸上了。}
}
{
    \eA{She Yüeh, on inspection, found indeed a hole burnt in it of the size of a
        finger. "This," she said, "must have been done by some spark from the
        hand-stove. It's of no consequence."}
}
{
}

\Verse{290}
{
    \zA{这不值什么，赶着叫人悄悄拿出去叫个能干织补匠人织上就是了。”}
}
{
    \eA{Immediately she called a servant to her. "Take this out on the sly," she
        bade her, "and let an experienced weaver patch it. It will be all right
        then." So saying, she packed it up in a wrapper,}
}
{
}

\Verse{291}
{
    \zA{说着，就用包 袱包了，叫了一个嬷嬷送出去，说：“赶天亮就有才好，千万别给老太太、太太知 道。”}
}
{
    \eA{and a nurse carried it outside. "It should be ready by daybreak," she urged.
        "And by no means let our old lady or Madame Wang know anything about it."
        The matron brought it back again, after a protracted absence. "Not only,"
        she explained;}
}
{
}

\Verse{292}
{
    \zA{婆子去了半日，仍就拿回来，说：“不但织补匠，能干裁缝、绣匠并做女工 的，问了，都不认的这是什么，都不敢揽。”}
}
{
    \eA{"have weavers, first-class tailors, and embroiderers, but even those, who do
        women's work, been asked about it, and they all have no idea what this is
        made of. None of them therefore will venture to undertake the job." "What's
        to be done?" She Yüeh inquired.}
}
{
}

\Verse{293}
{
    \zA{麝月道：“这怎么好呢？}
}
{
    \eA{"But it won't matter if you don't wear it to-morrow."}
}
{
}

\Verse{294}
{
    \zA{明儿不穿也 罢了。”}
}
{
    \eA{"To-morrow is the very day of the anniversary,"}
}
{
}

\Verse{295}
{
    \zA{宝玉道：“明儿是正日子，老太太、太太说了，还叫穿过这个去呢。}
}
{
    \eA{Pao-yü rejoined. "Grandmother and my mother bade me put this on and go and
        pay my visit; and here I go and burn it, on the first day I wear it. Now
        isn't this enough to throw a damper over my good cheer?"}
}
{
}

\Verse{296}
{
    \zA{偏头 一日就烧了，岂不扫兴！”}
}
{
    \eA{Ch'ing Wen lent an ear to their conversation for a long time, until unable
        to restrain herself,}
}
{
}

\Verse{297}
{
    \zA{晴雯听了半日，忍不住，翻身说道：“拿来我瞧瞧罢!没那福气穿就罢了，这 会子又着急。”}
}
{
    \eA{she twisted herself round. "Bring it here," she chimed in, "and let me see
        it! You haven't been lucky in wearing this; but never mind!" These words
        were still on Ch'ing Wen's lips, when the coat was handed to her.}
}
{
}

\Verse{298}
{
    \zA{宝玉笑道：“这话倒说的是。”}
}
{
    \eA{The lamp was likewise moved nearer to her. With minute care she surveyed it.}
}
{
}

\Verse{299}
{
    \zA{说着，便递给晴雯，又移过灯来， 细瞧了一瞧。}
}
{
    \eA{"This is made," Ch'ing Wen observed, "of gold thread, spun from peacock's
        feathers. So were we now to also take gold thread,}
}
{
}

\Verse{300}
{
    \zA{晴雯道：“这是孔雀金线的。}
}
{
    \eA{twisted from the feathers of the peacock, and darn it closely, by imitating
        the woof,}
}
{
}

\Verse{301}
{
    \zA{如今咱们也拿孔雀金线，就像界线似的 界密了，只怕还可混的过去。”}
}
{
    \eA{I think it will pass without detection." "The peacock-feather-thread is
        ready at hand," She Yüeh remarked smilingly. "But who's there, exclusive of
        you, able to join the threads?"}
}
{
}

\Verse{302}
{
    \zA{麝月笑道：“孔雀线现成的，但这里除你，还有谁 会界线？”}
}
{
    \eA{"I'll, needless to say, do my level best to the very cost of my life and
        finish," Ch'ing Wen added. "How ever could this do?" Pao-yü eagerly
        interposed. "You're just slightly better,}
}
{
}

\Verse{303}
{
    \zA{晴雯道：“说不的我挣命罢了。”}
}
{
    \eA{and how could you take up any needlework?" "You needn't go on in this
        chicken-hearted way!"}
}
{
}

\Verse{304}
{
    \zA{宝玉忙道：“这如何使得？}
}
{
    \eA{Ch'ing Wen cried. "I know my own self well enough." With this reply, she sat
        up,}
}
{
}

\Verse{305}
{
    \zA{才好了些， 如何做得活！”}
}
{
    \eA{and, putting her hair up, she threw something over her shoulders.}
}
{
}

\Verse{306}
{
    \zA{晴雯道：“不用你蝎蝎螫螫的，我自知道。”}
}
{
    \eA{Her head felt heavy; her body light. Before her eyes, confusedly flitted
        golden stirs. In real deed, she could not stand the strain.}
}
{
}

\Verse{307}
{
    \zA{一面说，一面坐起来， 挽了一挽头发，披了衣裳。}
}
{
    \eA{But when inclined to give up the work, she again dreaded that Pao-yü would
        be driven to despair. She therefore had perforce to make a supreme effort
        and,}
}
{
}

\Verse{308}
{
    \zA{只觉头重身轻，满眼金星乱迸，实实掌不住。}
}
{
    \eA{setting her teeth to, she bore the exertion. All the help she asked of She
        Yüeh was to lend her a hand in reeling the thread.}
}
{
}

\Verse{309}
{
    \zA{待不做， 又怕宝玉着急，少不得狠命咬牙捱着。}
}
{
    \eA{Ch'ing Wen first took hold of a thread, and put it side by side (with those
        in the pelisse) to compare the two together. "This," she remarked, "isn't
        quite like them;}
}
{
}

\Verse{310}
{
    \zA{便命麝月只帮着拈线。}
}
{
    \eA{but when it's patched up with it, it won't show very much."}
}
{
}

\Verse{311}
{
    \zA{晴雯先拿了一根比一 比，笑道：“这虽不很像，要补上也不很显。”}
}
{
    \eA{"It will do very well," Pao-yü said. "Could one also go and hunt up a
        Russian tailor?" Ch'ing Wen commenced by unstitching the lining, and,
        inserting under it,}
}
{
}

\Verse{312}
{
    \zA{宝玉道：“这就很好，那里又找俄 罗斯国的裁缝去？”}
}
{
    \eA{a bamboo bow, of the size of the mouth of a tea cup, she bound it tight at
        the back. She then turned her mind to the four sides of the aperture,}
}
{
}

\Verse{313}
{
    \zA{晴雯先将里子拆开，用茶杯口大小一个竹弓钉绷在背面，再将 破口四边用金刀刮的散松松的，然后用针缝了两条，分出经纬，亦如界线之法，先
        界出地子来，后依本纹来回织补。}
}
{
    \eA{and these she loosened by scratching them with a golden knife. Making next
        two stitches across with her needle, she marked out the warp and woof; and,
        following the way the threads were joined, she first and foremost connected
        the foundation,}
}
{
}

\Verse{314}
{
    \zA{补两针，又看看；织补不上三五针，便伏在枕上 歇一会。}
}
{
    \eA{and then keeping to the original lines, she went backwards and forwards
        mending the hole; passing her work, after every second stitch,}
}
{
}

\Verse{315}
{
    \zA{宝玉在旁，一时又问：“吃些滚水不吃？”}
}
{
    \eA{under further review. But she did not ply her needle three to five times,
        before she lay herself down on her pillow,}
}
{
}

\Verse{316}
{
    \zA{一时又命：“歇一歇。”}
}
{
    \eA{and indulged in a little rest. Pao-yü was standing by her side.}
}
{
}

\Verse{317}
{
    \zA{一时 又拿一件灰鼠斗篷替他披在背上，一时又拿个枕头给他靠着。}
}
{
    \eA{Now he inquired of her: "Whether she would like a little hot water to
        drink." Later on, he asked her to repose herself. Now he seized a
        grey-squirrel wrapper and threw it over her shoulders.}
}
{
}

\Verse{318}
{
    \zA{急的晴雯央道：“小 祖宗，你只管睡罢!再熬上半夜，明儿眼睛抠搂了，那恰怎么好？”}
}
{
    \eA{Shortly after, he took a pillow and propped her up. (The way he fussed) so
        exasperated Ch'ing Wen that she begged and entreated him to leave off. "My
        junior ancestor!" she exclaimed, "do go to bed and sleep! If you sit up for
        the other half of the night,}
}
{
}

\Verse{319}
{
    \zA{宝玉见他着急，只得胡乱睡下，仍睡不着。}
}
{
    \eA{your eyes will to-morrow look as if they had been scooped out, and what good
        will possibly come out of that?"}
}
{
}

\Verse{320}
{
    \zA{一时只听自鸣钟已敲了四下，刚刚 补完；又用小牙刷慢慢的剔出？}
}
{
    \eA{Pao-yü realised her state of exasperation and felt compelled to come and lie
        down anyhow. But he could not again close his eyes. In a little while, she
        heard the clock strike four,}
}
{
}

\Verse{321}
{
    \zA{毛来。}
}
{
    \eA{and just managing to finish she took a small tooth-brush,}
}
{
}

\Verse{322}
{
    \zA{麝月道：“这就很好，要不留心，再看不出 的。”}
}
{
    \eA{and rubbed up the pile. "That will do!" She Yüeh put in. "One couldn't
        detect it, unless one examined it carefully."}
}
{
}

\Verse{323}
{
    \zA{宝玉忙要了瞧瞧，笑说：“真真一样了。”}
}
{
    \eA{Pao-yü asked with alacrity to be allowed to have a look at it. "Really," he
        smiled, "it's quite the same thing."}
}
{
}

\Verse{324}
{
    \zA{晴雯已嗽了几声，好容易补完了， 说了一声：“补虽补了，到底不像。}
}
{
    \eA{Ch'ing Wen coughed and coughed time after time, so it was only after extreme
        difficulty that she succeeded in completing what she had to patch. "It's
        mended, it's true,"}
}
{
}

\Verse{325}
{
    \zA{我也再不能了！”}
}
{
    \eA{she remarked, "but it does not, after all, look anything like it.}
}
{
}

\Verse{326}
{
    \zA{“嗳哟”了一声，就身不由 主睡下了。}
}
{
    \eA{Yet, I cannot stand the effort any more!" As she shouted 'Ai-ya,' she lost
        control over herself,}
}
{
}

\Verse{327}
{
    \zA{要知端的，且看下回分解。}
}
{
    \eA{and dropped down upon the bed. But, reader, if you choose to know anything
        more of her state,}
}
{
}

\Verse{328}
{
    \zA{(本章完)}
}
{
    \eA{peruse the next chapter.}
}
{
}
